\ExerciseSection

\begin{enumerate}
\def\labelenumi{\arabic{enumi})}
\tightlist
\item
  设 X 是一个集，Y 是一个含多于一个点的一致空间，S 是 X 的非空子集所成的一个非空集。设 \(Y' \subset \mathcal{F}(X; Y)\) 是 X 到 Y 中的常值映射全体所成之集。对每个 \(y \in Y\) ，设 \(c_{y}\) 是以 y 为值的 X 到 Y 中的常值映射。
\end{enumerate}

\begin{enumerate}
\def\labelenumi{\alph{enumi})}
\item
  证明 \(y \to c_{y}\) 是 Y 到一致子空间 \(Y'\) ( \(\mathcal{F}_{\mathfrak{S}}(\mathbf{X}; \mathbf{Y})\) 的一致子空间)上的一个同构。
\item
  \(\mathcal{F}_{\mathfrak{S}}(\mathbf{X};\mathbf{Y})\) 是豪斯多夫的，当且仅当 Y 是豪斯多夫的且 \(\mathfrak{S}\) 是 X 的一个覆盖。
\item
  \(\mathbf{Y}'\) 在 \(\mathcal{F}_{\mathfrak{S}}(\mathbf{X};\mathbf{Y})\) 中是闭的，当且仅当 \(\mathcal{F}_{\mathfrak{S}}(\mathbf{X};\mathbf{Y})\) 是豪斯多夫的。
\end{enumerate}

\begin{enumerate}
\def\labelenumi{\arabic{enumi})}
\setcounter{enumi}{1}
\tightlist
\item
  设 X 是一个集，Y 是一个含多于一个点的豪斯多夫一致空间，\(\mathfrak{S}_1, \mathfrak{S}_2\) 是 X 的子集所成的两个集，它们满足第 2 号的条件 \((\mathrm{F}_{\mathrm{I}}^{\prime}), (\mathrm{F}_{\mathrm{II}}^{\prime})\) 。证明若 \(\mathfrak{S}_1 \subset \mathfrak{S}_2\) 且 \(\mathfrak{S}_1 \neq \mathfrak{S}_2\) ，则 \(\mathfrak{S}_1\) -收敛的一致结构严格粗于 \(\mathfrak{S}_2\) -收敛的一致结构。特别地：
\end{enumerate}

\begin{enumerate}
\def\labelenumi{(\roman{enumi})}
\item
  若 X 是非紧的豪斯多夫拓扑空间，则紧收敛的一致结构严格粗于一致收敛的一致结构。
\item
  若 X 是具有无限紧子集的豪斯多夫拓扑空间（参见第 I 章，§ 9，习题 4），则逐点收敛的一致结构严格粗于紧收敛的一致结构。
\end{enumerate}

\begin{enumerate}
\def\labelenumi{\arabic{enumi})}
\setcounter{enumi}{2}
\item
  设 X 是一个集，S 是 X 的一个覆盖，Y 是一个非豪斯多夫的一致空间，\(Y_{0}\) 是与 Y 相伴的豪斯多夫一致空间（第 II 章，§ 3，第 8 号）。证明与 \(\mathcal{F}_{\mathfrak{S}}(\mathrm{X};\mathrm{Y})\) 相伴的豪斯多夫一致空间同构于 \(\mathcal{F}_{\mathfrak{S}}(\mathrm{X};\mathrm{Y}_{0})\) 。
\item
  证明在集 \(\mathcal{C}(\mathbf{R};\mathbf{R})\) (即定义在 \(\mathbf{R}\) 上的有限连续实值函数全体所成之集)上，关于 R 的加法一致结构的一致收敛拓扑不同于关于由扩充实直线 \(\overline{R}\) 的（唯一的）一致结构在 R 上诱导的一致结构的一致收敛拓扑（尽管这两个一致结构在 R 上诱导的拓扑相同）。
\item
  设 X 是拓扑空间。X 的子集所成的集 S 称为饱和的，如果它满足第 2 号的条件 \((\mathbf{F}_{\mathbf{I}}^{\prime})\) 与 \((\mathbf{F}_{\mathbf{II}}^{\prime})\) ，并且 S 中每个集的闭包属于 S。
\end{enumerate}

\begin{enumerate}
\def\labelenumi{\alph{enumi})}
\item
  证明若 X 是正规的且 S 是饱和的并覆盖 X，则集 C(X; R) 在空间 \(\tilde{\mathcal{C}}_{\mathfrak{S}}(X; \mathbf{R})\) (第 6 号，定理 2，推论 2)中关于 S-收敛拓扑是稠密的。特别地，C(X; R) 在 \(F_{s}(X; \mathbf{R})\) 中稠密。只有当 X 是离散的时，C(X; R) 才在 \(F_{s}(X; \mathbf{R})\) 中是闭的。
\item
  设 X 是完全正则空间（第 IX 章，§ 1，第 5 号），\(\mathfrak{S}_1, \mathfrak{S}_2\) 是 X 的子集所成的两个饱和集，满足 \(\mathfrak{S}_1 \subset \mathfrak{S}_2\) 且 \(\mathfrak{S}_1 \neq \mathfrak{S}_2\) 。证明 \(\mathfrak{S}_1\) -收敛在 \(\mathcal{C}(\mathbf{X}; \mathbf{R})\) 上的拓扑严格粗于 \(\mathfrak{S}_2\) -收敛的拓扑。
\end{enumerate}

\begin{enumerate}
\def\labelenumi{\arabic{enumi})}
\setcounter{enumi}{5}
\tightlist
\item
  \begin{enumerate}
  \def\labelenumii{\alph{enumii})}
  \tightlist
  \item
    设 X 是拓扑空间，Y 是一致空间，\(\Phi\) 是 \(\mathcal{C}(X;Y)\) 所成之集上逐点收敛于函数 \(u_{0}\) 的一个滤子。则 \(u_{0}\) 在一点 \(x_{0}\in X\) 处连续，当且仅当对 Y 的每个围域 V 与每个集 \(M\in\Phi\) ，存在 \(x_{0}\) 的一个邻域 U 与一个映射 \(u\in M\) ，使得 \((u_{0}(x),u(x))\in V\) 对一切 \(x\in U\) 成立。
  \end{enumerate}
\end{enumerate}

\begin{enumerate}
\def\labelenumi{\alph{enumi})}
\setcounter{enumi}{1}
\tightlist
\item
  设 X 是拟紧空间，Y 是一致空间，\(\Phi\) 是 \(\mathcal{C}(X;Y)\) 上逐点收敛于函数 \(u_{0}\) 的一个滤子。则 \(u_{0}\) 在 X 上连续，当且仅当对 Y 的每个围域 V 与每个集 \(M\in\Phi\) ，存在有限个函数 \(u_{i}\in\mathbf{M}(1\leqslant i\leqslant n)\) ，使得对每个 \(x\in X\) ，至少对一个指标 i 有 \((u_{0}(x),u_{i}(x))\in V\) 。
\end{enumerate}

\begin{enumerate}
\def\labelenumi{\arabic{enumi})}
\setcounter{enumi}{6}
\tightlist
\item
  设 X、Y 是两个豪斯多夫一致空间。对每个连续映射 \(f: \mathbf{X} \to \mathbf{Y}\) ，设 \(\mathbf{G}(f) \subset \mathbf{X} \times \mathbf{Y}\) 是 \(f\) 的图象；\(\mathbf{G}(f)\) 在 \(\mathbf{X} \times \mathbf{Y}\) 中是闭的（第 I 章，§ 8，第 1 号，命题 2，推论 2）：因而 \(f \to \mathbf{G}(f)\) 是 \(\mathcal{C}(\mathbf{X}; \mathbf{Y})\) 到 \(\mathfrak{F}(\mathbf{X} \times \mathbf{Y})\) (即 \(\mathbf{X} \times \mathbf{Y}\) 的非空闭子集全体所成之集)中的一个单射映射。
\end{enumerate}

\begin{enumerate}
\def\labelenumi{\alph{enumi})}
\item
  证明映射 \(f \to \mathbf{G}(f)\) (它把 \(\mathcal{C}_u(\mathbf{X};\mathbf{Y})\) 映到 \(\mathfrak{F}(\mathbf{X}\times \mathbf{Y})\) 中)当 \(\mathfrak{F}(\mathbf{X}\times \mathbf{Y})\) 赋有第 II 章，§ 1，习题 5 中定义的一致结构时是一致连续的。
\item
  设 \(\Gamma\) 是 \(\mathcal{C}(\mathbf{X};\mathbf{Y})\) 在 \(\mathfrak{F}(\mathbf{X}\times \mathbf{Y})\) 中 (在映射 \(f\to \mathbf{G}(f)\) 下)的象，并设 \(\varphi :\Gamma \to \mathcal{C}_u(\mathbf{X};\mathbf{Y})\) 是逆映射。证明若 \(\mathbf{X}\) 紧，则 \(\varphi\) 在 \(\Gamma\) 上连续（用反证法论证）。
\item
  取 X 与 Y 都为 R 的紧区间 {[}0, 1{]} 。证明 \(\varphi\) 在 \(\Gamma\) 上不是一致连续的。
\end{enumerate}

\begin{enumerate}
\def\labelenumi{\arabic{enumi})}
\setcounter{enumi}{7}
\tightlist
\item
  设 X、Y 是两个度量空间，\((f_{n})\) 是 X 到 Y 中的 \(\alpha\) 类波莱尔映射的一个序列（第 IX 章，§ 6，习题 16）。
\end{enumerate}

\begin{enumerate}
\def\labelenumi{\alph{enumi})}
\item
  假设序列 \((f_{n})\) 逐点收敛于映射 f: X → Y。证明 f 是 \(\alpha + 1\) 类的。{[}若 U 是 Y 的开子集，注意 \(\overline{f}^{1}(U) = \bigcup_{n \geqslant 1} \left( \bigcap_{k \geqslant 0} \overline{f}_{n+k}^{1}(U) \right)\) ，并利用第 IX 章，§ 6，习题 4。{]}
\item
  假设序列 \((f_n)\) 一致收敛于 \(f\) 。证明 \(f\) 是 \(\alpha\) 类的。{[}设 \(\mathbf{F}\) 是 \(\mathbf{Y}\) 的一个闭子集，且对每个整数 \(n > 0\) ，设 \(\mathbf{V}_n\) 是 \(\mathbf{Y}\) 中与 \(\mathbf{F}\) 的距离 \(\leqslant 1/n\) 的点所成之集。证明存在整数的一个递增序列 \(n \to m(n)\) 使得 \(\overline{f}^1(\mathbf{F}) = \bigcap_{n} \overline{f}_{m(n)}^1(\mathbf{V}_n)\) 。{]}
\item
  假设 Y 是可数型的。证明若 \(f: X \to Y\) 是 \(\alpha > 0\) 类的波莱尔函数，则存在波莱尔函数的一个序列 \(g_{n}: X \to Y\) ，它们是 \(<\alpha\) 类的，使得 \((g_{n})\) 逐点收敛于 f。{[}证明可取 \(g_{n}\) 为只取有限多个值的函数；利用第 IX 章，§ 6 的习题 4 h) 与 16 c)。{]}
\end{enumerate}

\begin{enumerate}
\def\labelenumi{\arabic{enumi})}
\setcounter{enumi}{8}
\item
  设 X 是贝尔空间，Y 是度量空间，\((f_{n})\) 是 X 到 Y 中的连续映射的一个序列，它逐点收敛于函数 f。点 \(x \in X\) 称为序列 \((f_{n})\) 的一致收敛点，如果对每个 \(\varepsilon > 0\) ，存在 x 的一个邻域 V 与一个整数 p，使得 \(d(f_{m}(y), f_{n}(y)) \leqslant \varepsilon\) 对一切 \(y \in V\) 与一切整数 \(m \geqslant p, n \geqslant p, d\) 成立，其中 d 是 Y 上的度量。证明序列 \((f_{n})\) 的一致收敛点所成之集的余集 S 在 X 中是贫集{[}参见第 IX 章，§ 5，习题 22 a{]}。给出一个 S 在 X 中稠密的例子。{[}取 X = Y = R，设 \(n \to r_{n}\) 是 N 到 Q 上的一个双射，并设 \((g_{n})\) 是 R 到 {[}0, 1{]} 中的连续映射的一个序列，它对 \(x \neq 0\) 收敛于零函数而在 x = 0 处收敛于 1，且在每个不含 0 的 R 的开集中收敛是一致的。然后考察函数序列 \(f_{n}(x) = \sum_{p=0}^{\infty} \alpha_{p} g_{n}(x - r_{p})\) ，其中序列 \((\alpha_{p})\) 以适当的方式趋于 0。{]}
\item
  证明在群 \(\Gamma\) (由实直线 \(\mathbf{R}\) 到其自身上的同胚组成)上，逐点收敛拓扑与紧收敛拓扑相同（参见 § 3，习题 14）。
\item
  设 X 是拓扑空间，G 是拓扑群。X 到 G 中的连续映射所成之集 C(X; G) 是群 \(G^{X}\) 的一个子群。设 S 是 X 的子集所成的一个集。
\end{enumerate}

\begin{enumerate}
\def\labelenumi{\alph{enumi})}
\item
  假设对每个 \(A \in S\) 、G 中 e 的每个邻域 V 与每个 \(u \in C(X; G)\) ，存在 G 中 e 的一个邻域 W，使得 \(sWs^{-1} \subset W\) 对一切 \(s \in u(A)\) 成立。证明 S-收敛拓扑这时与 \(\mathcal{C}(X; G)\) 的群结构相容，且如此定义的拓扑群 \(\mathcal{C}_{\mathfrak{S}}(X; G)\) 的右（相应地，左）一致结构与 S-收敛的一致结构相同。考察逐点收敛的情形、G 局部紧时紧收敛的情形，以及 G 交换的情形。
\item
  若 G 是群 \(\mathbf{SL}(2, \mathbf{R})\) ，赋有由 \(R^{4}\) 的拓扑诱导的拓扑，证明一致收敛拓扑与 \(\mathcal{C}(\mathbf{R}; \mathbf{G})\) 的群结构不相容。
\end{enumerate}

\begin{enumerate}
\def\labelenumi{\arabic{enumi})}
\setcounter{enumi}{11}
\tightlist
\item
  设 X 是拓扑空间，A 是拓扑环（第 III 章，§ 6，第 3 号）。X 到 A 中的连续映射所成之集 C(X; A) 是环 A\^{}X 的一个子环。设 S 是 X 的子集所成的一个集。
\end{enumerate}

\begin{enumerate}
\def\labelenumi{\alph{enumi})}
\item
  假设对每个 \(\mathbf{M} \in \mathfrak{S}\) 与每个 \(u \in \mathcal{C}(\mathbf{X}; \mathbf{A})\) ，集 \(u(\mathbf{M})\) 是有界的（第 III 章，§ 6，习题 12）。证明 \(\mathfrak{S}\) -收敛拓扑与 \(\mathcal{C}(\mathbf{X}; \mathbf{A})\) 的环结构相容。考察逐点收敛的情形与紧收敛的情形。
\item
  设 X 是一个非紧的局部紧 \(\sigma\) -紧空间（第 I 章，§ 9，第 9 号）。证明一致收敛拓扑与 \(\mathcal{C}\) (X; R) 的环结构不相容。
\end{enumerate}

